\IfFileExists{stacks-project.cls}{%
\documentclass{stacks-project}
}{%
\documentclass{amsart}
}

% For dealing with references we use the comment environment
\usepackage{verbatim}
\newenvironment{reference}{\comment}{\endcomment}
%\newenvironment{reference}{}{}
\newenvironment{slogan}{\comment}{\endcomment}
\newenvironment{history}{\comment}{\endcomment}

% For commutative diagrams we use Xy-pic
\usepackage[all]{xy}

% We use 2cell for 2-commutative diagrams.
\xyoption{2cell}
\UseAllTwocells

% We use multicol for the list of chapters between chapters
\usepackage{multicol}

% This is generally recommended for better output
\usepackage{lmodern}
\usepackage[T1]{fontenc}

% For cross-file-references
\usepackage{xr-hyper}

% Package for hypertext links:
\usepackage{hyperref}

% For any local file, say "hello.tex" you want to link to please
% use \externaldocument[hello-]{hello}
\externaldocument[introduction-]{introduction}
\externaldocument[conventions-]{conventions}
\externaldocument[sets-]{sets}
\externaldocument[categories-]{categories}
\externaldocument[topology-]{topology}
\externaldocument[sheaves-]{sheaves}
\externaldocument[sites-]{sites}
\externaldocument[stacks-]{stacks}
\externaldocument[fields-]{fields}
\externaldocument[algebra-]{algebra}
\externaldocument[brauer-]{brauer}
\externaldocument[homology-]{homology}
\externaldocument[derived-]{derived}
\externaldocument[simplicial-]{simplicial}
\externaldocument[more-algebra-]{more-algebra}
\externaldocument[smoothing-]{smoothing}
\externaldocument[modules-]{modules}
\externaldocument[sites-modules-]{sites-modules}
\externaldocument[injectives-]{injectives}
\externaldocument[cohomology-]{cohomology}
\externaldocument[sites-cohomology-]{sites-cohomology}
\externaldocument[dga-]{dga}
\externaldocument[dpa-]{dpa}
\externaldocument[sdga-]{sdga}
\externaldocument[hypercovering-]{hypercovering}
\externaldocument[schemes-]{schemes}
\externaldocument[constructions-]{constructions}
\externaldocument[properties-]{properties}
\externaldocument[morphisms-]{morphisms}
\externaldocument[coherent-]{coherent}
\externaldocument[divisors-]{divisors}
\externaldocument[limits-]{limits}
\externaldocument[varieties-]{varieties}
\externaldocument[topologies-]{topologies}
\externaldocument[descent-]{descent}
\externaldocument[perfect-]{perfect}
\externaldocument[more-morphisms-]{more-morphisms}
\externaldocument[flat-]{flat}
\externaldocument[groupoids-]{groupoids}
\externaldocument[more-groupoids-]{more-groupoids}
\externaldocument[etale-]{etale}
\externaldocument[chow-]{chow}
\externaldocument[intersection-]{intersection}
\externaldocument[pic-]{pic}
\externaldocument[weil-]{weil}
\externaldocument[adequate-]{adequate}
\externaldocument[dualizing-]{dualizing}
\externaldocument[duality-]{duality}
\externaldocument[discriminant-]{discriminant}
\externaldocument[derham-]{derham}
\externaldocument[local-cohomology-]{local-cohomology}
\externaldocument[algebraization-]{algebraization}
\externaldocument[curves-]{curves}
\externaldocument[resolve-]{resolve}
\externaldocument[models-]{models}
\externaldocument[functors-]{functors}
\externaldocument[equiv-]{equiv}
\externaldocument[pione-]{pione}
\externaldocument[etale-cohomology-]{etale-cohomology}
\externaldocument[proetale-]{proetale}
\externaldocument[relative-cycles-]{relative-cycles}
\externaldocument[more-etale-]{more-etale}
\externaldocument[trace-]{trace}
\externaldocument[crystalline-]{crystalline}
\externaldocument[spaces-]{spaces}
\externaldocument[spaces-properties-]{spaces-properties}
\externaldocument[spaces-morphisms-]{spaces-morphisms}
\externaldocument[decent-spaces-]{decent-spaces}
\externaldocument[spaces-cohomology-]{spaces-cohomology}
\externaldocument[spaces-limits-]{spaces-limits}
\externaldocument[spaces-divisors-]{spaces-divisors}
\externaldocument[spaces-over-fields-]{spaces-over-fields}
\externaldocument[spaces-topologies-]{spaces-topologies}
\externaldocument[spaces-descent-]{spaces-descent}
\externaldocument[spaces-perfect-]{spaces-perfect}
\externaldocument[spaces-more-morphisms-]{spaces-more-morphisms}
\externaldocument[spaces-flat-]{spaces-flat}
\externaldocument[spaces-groupoids-]{spaces-groupoids}
\externaldocument[spaces-more-groupoids-]{spaces-more-groupoids}
\externaldocument[bootstrap-]{bootstrap}
\externaldocument[spaces-pushouts-]{spaces-pushouts}
\externaldocument[spaces-chow-]{spaces-chow}
\externaldocument[groupoids-quotients-]{groupoids-quotients}
\externaldocument[spaces-more-cohomology-]{spaces-more-cohomology}
\externaldocument[spaces-simplicial-]{spaces-simplicial}
\externaldocument[spaces-duality-]{spaces-duality}
\externaldocument[formal-spaces-]{formal-spaces}
\externaldocument[restricted-]{restricted}
\externaldocument[spaces-resolve-]{spaces-resolve}
\externaldocument[formal-defos-]{formal-defos}
\externaldocument[defos-]{defos}
\externaldocument[cotangent-]{cotangent}
\externaldocument[examples-defos-]{examples-defos}
\externaldocument[algebraic-]{algebraic}
\externaldocument[examples-stacks-]{examples-stacks}
\externaldocument[stacks-sheaves-]{stacks-sheaves}
\externaldocument[criteria-]{criteria}
\externaldocument[artin-]{artin}
\externaldocument[quot-]{quot}
\externaldocument[stacks-properties-]{stacks-properties}
\externaldocument[stacks-morphisms-]{stacks-morphisms}
\externaldocument[stacks-limits-]{stacks-limits}
\externaldocument[stacks-cohomology-]{stacks-cohomology}
\externaldocument[stacks-perfect-]{stacks-perfect}
\externaldocument[stacks-introduction-]{stacks-introduction}
\externaldocument[stacks-more-morphisms-]{stacks-more-morphisms}
\externaldocument[stacks-geometry-]{stacks-geometry}
\externaldocument[moduli-]{moduli}
\externaldocument[moduli-curves-]{moduli-curves}
\externaldocument[examples-]{examples}
\externaldocument[exercises-]{exercises}
\externaldocument[guide-]{guide}
\externaldocument[desirables-]{desirables}
\externaldocument[coding-]{coding}
\externaldocument[obsolete-]{obsolete}
\externaldocument[fdl-]{fdl}
\externaldocument[index-]{index}

% Theorem environments.
%
\theoremstyle{plain}
\newtheorem{theorem}[subsection]{Theorem}
\newtheorem{proposition}[subsection]{Proposition}
\newtheorem{lemma}[subsection]{Lemma}

\theoremstyle{definition}
\newtheorem{definition}[subsection]{Definition}
\newtheorem{example}[subsection]{Example}
\newtheorem{exercise}[subsection]{Exercise}
\newtheorem{situation}[subsection]{Situation}

\theoremstyle{remark}
\newtheorem{remark}[subsection]{Remark}
\newtheorem{remarks}[subsection]{Remarks}

\numberwithin{equation}{subsection}

% Macros
%
\def\lim{\mathop{\mathrm{lim}}\nolimits}
\def\colim{\mathop{\mathrm{colim}}\nolimits}
\def\Spec{\mathop{\mathrm{Spec}}}
\def\Hom{\mathop{\mathrm{Hom}}\nolimits}
\def\Ext{\mathop{\mathrm{Ext}}\nolimits}
\def\SheafHom{\mathop{\mathcal{H}\!\mathit{om}}\nolimits}
\def\SheafExt{\mathop{\mathcal{E}\!\mathit{xt}}\nolimits}
\def\Sch{\mathit{Sch}}
\def\Mor{\mathop{\mathrm{Mor}}\nolimits}
\def\Ob{\mathop{\mathrm{Ob}}\nolimits}
\def\Sh{\mathop{\mathit{Sh}}\nolimits}
\def\NL{\mathop{N\!L}\nolimits}
\def\CH{\mathop{\mathrm{CH}}\nolimits}
\def\proetale{{pro\text{-}\acute{e}tale}}
\def\etale{{\acute{e}tale}}
\def\QCoh{\mathit{QCoh}}
\def\Ker{\mathop{\mathrm{Ker}}}
\def\Im{\mathop{\mathrm{Im}}}
\def\Coker{\mathop{\mathrm{Coker}}}
\def\Coim{\mathop{\mathrm{Coim}}}

% Boxtimes
%
\DeclareMathSymbol{\boxtimes}{\mathbin}{AMSa}{"02}

%
% Macros for moduli stacks/spaces
%
\def\QCohstack{\mathcal{QC}\!\mathit{oh}}
\def\Cohstack{\mathcal{C}\!\mathit{oh}}
\def\Spacesstack{\mathcal{S}\!\mathit{paces}}
\def\Quotfunctor{\mathrm{Quot}}
\def\Hilbfunctor{\mathrm{Hilb}}
\def\Curvesstack{\mathcal{C}\!\mathit{urves}}
\def\Polarizedstack{\mathcal{P}\!\mathit{olarized}}
\def\Complexesstack{\mathcal{C}\!\mathit{omplexes}}
% \Pic is the operator that assigns to X its picard group, usage \Pic(X)
% \Picardstack_{X/B} denotes the Picard stack of X over B
% \Picardfunctor_{X/B} denotes the Picard functor of X over B
\def\Pic{\mathop{\mathrm{Pic}}\nolimits}
\def\Picardstack{\mathcal{P}\!\mathit{ic}}
\def\Picardfunctor{\mathrm{Pic}}
\def\Deformationcategory{\mathcal{D}\!\mathit{ef}}
\usepackage{cleveref}
\crefname{originaldefinition}{original definition}{original definitions}
\Crefname{originaldefinition}{Original definition}{Original definitions}
\makeatletter
\newlabel{more-morphisms-definition-weighting}{{0F3A}{source}{}{sourcepointer-1}{}}
\newlabel{more-morphisms-definition-weighting@cref}{{[originaldefinition][0][]0F3A}{[1][0][]source}{}{}{}}
\makeatother
\crefname{originallemma}{original lemma}{original lemmas}
\Crefname{originallemma}{Original lemma}{Original lemmas}
\makeatletter
\newlabel{lemma-weightings-pre}{{lemma-weightings-pre}{source}{}{sourcepointer-2}{}}
\newlabel{lemma-weightings-pre@cref}{{[originallemma][0][]lemma-weightings-pre}{[1][0][]source}{}{}{}}
\makeatother
\crefname{originallemma}{original lemma}{original lemmas}
\Crefname{originallemma}{Original lemma}{Original lemmas}
\makeatletter
\newlabel{lemma-relative-cycle-smooth}{{0H5B}{source}{}{sourcepointer-3}{}}
\newlabel{lemma-relative-cycle-smooth@cref}{{[originallemma][0][]0H5B}{[1][0][]source}{}{}{}}
\makeatother
\crefname{originallemma}{original lemma}{original lemmas}
\Crefname{originallemma}{Original lemma}{Original lemmas}
\makeatletter
\newlabel{morphisms-lemma-locally-quasi-finite-qc-source-universally-bounded}{{03JA}{source}{}{sourcepointer-4}{}}
\newlabel{morphisms-lemma-locally-quasi-finite-qc-source-universally-bounded@cref}{{[originallemma][0][]03JA}{[1][0][]source}{}{}{}}
\makeatother
\crefname{originallemma}{original lemma}{original lemmas}
\Crefname{originallemma}{Original lemma}{Original lemmas}
\makeatletter
\newlabel{more-morphisms-lemma-weighting-base-change}{{0F3B}{source}{}{sourcepointer-5}{}}
\newlabel{more-morphisms-lemma-weighting-base-change@cref}{{[originallemma][0][]0F3B}{[1][0][]source}{}{}{}}
\makeatother
\crefname{originallemma}{original lemma}{original lemmas}
\Crefname{originallemma}{Original lemma}{Original lemmas}
\makeatletter
\newlabel{more-morphisms-lemma-weighting-specialization}{{more-morphisms-lemma-weighting-specialization}{source}{}{sourcepointer-6}{}}
\newlabel{more-morphisms-lemma-weighting-specialization@cref}{{[originallemma][0][]more-morphisms-lemma-weighting-specialization}{[1][0][]source}{}{}{}}
\makeatother
\title{Weightings of locally quasi-finite morphisms and relative zero-cycles}
\author{The Stacks Project Authors}
\date{stacks-project:relative-cycles.tex:local-snapshot-20261012; GNU FDL 1.2 or later}
\begin{document}
\maketitle
\setcounter{subsection}{1}
\noindent\textit{Editorial note.} All delivered mathematical text is the author's own, copied from the pinned Stacks Project chapter file relative-cycles.tex; no statement or proof step was rewritten, reordered or removed. The counted span is the author environment and its own proof environment. References inside the span are handed to the assembler as navigation stubs that carry the canonical tag of the cited tag only; whether each of them is load-bearing is recorded on the stub itself.\par\medskip
\noindent\textit{Allowed standard prerequisites (editorial).}\begin{itemize}
\item Weighting definition and diagram: for a locally quasi-finite morphism a weighting is tested on an etale localization diagram with pi : U to V (More on Morphisms, Definition more-morphisms-definition-weighting, canonical tag 0F3A). cited at proof line 1744. Fixes the diagram in which the weighting w\_alpha must be shown to have locally constant integral over pi.
\item Base-change compatibility of the weighting of a relative zero-cycle: w\_\{g\textasciicircum{}*alpha\} = w\_alpha composed with g' for a morphism g : S' to S of locally Noetherian schemes (relative-cycles chapter, Lemma lemma-weightings-pre, no canonical tag in the snapshot tag table, defined at relative-cycles source line 1671). cited at proof lines 1746 and 1778. First use: reduces part (1) to local constancy on V. Second use: its proof never used that alpha is a relative cycle, so w\_alpha = n w persists after base change in part (2).
\item Local constancy along smooth morphisms: the support behaviour of relative cycles along a smooth morphism of relative dimension (relative-cycles chapter, Lemma lemma-relative-cycle-smooth, canonical tag 0H5B). cited at proof line 1758. The computed fibre sum is the coefficient of v in pi\_*beta, and the lemma makes that function locally constant on V.
\item Universal boundedness of fibres: a morphism locally quasi-finite with quasi-compact source has universally bounded fibres (Morphisms chapter, Lemma morphisms-lemma-locally-quasi-finite-qc-source-universally-bounded, canonical tag 03JA). cited at proof lines 1770-1773. Guarantees a sufficiently divisible integer n making the constructed zero-cycles integral.
\item Base change of a weighting is a weighting (More on Morphisms, Lemma more-morphisms-lemma-weighting-base-change, canonical tag 0F3B). cited at proof lines 1780-1782. Lets the reduction to the discrete valuation ring case keep the identity w\_alpha = n w after base change.
\item Specialization lemma for weightings: a weighting of a locally quasi-finite morphism vanishing at all points over a point y vanishes at all points over every specialization y' of y (More on Morphisms, Lemma more-morphisms-lemma-weighting-specialization, no canonical tag in the snapshot tag table, defined at more-morphisms source line 22469). cited at proof lines 1795-1797. Kills the difference weighting w' = w\_beta - n w, which vanishes over the generic point of the discrete valuation ring, forcing w\_beta = n w and beta = alpha.
\end{itemize}\medskip
\ifdefined\hypertarget\hypertarget{sourcepointer-1}{}\fi
\noindent\textit{Source reference.} Original definition 0F3A: sources/stacks-chow/tags.txt, source lines 18860--18860. Cited at proof line 1744: fixes the diagram in which the weighting is tested.\par\medskip
\ifdefined\hypertarget\hypertarget{sourcepointer-2}{}\fi
\noindent\textit{Source reference.} Original lemma lemma-weightings-pre: sources/stacks-relative-cycles/pinned-original-source.tex, source lines 1671--1671. Cited at proof lines 1746 and 1778: base-change compatibility of the weighting, and the fact its proof never used the relative-cycle property.\par\medskip
\ifdefined\hypertarget\hypertarget{sourcepointer-3}{}\fi
\noindent\textit{Source reference.} Original lemma 0H5B: sources/stacks-chow/tags.txt, source lines 21381--21381. Cited at proof line 1758: makes the coefficient function of the pushed-forward cycle locally constant on V.\par\medskip
\ifdefined\hypertarget\hypertarget{sourcepointer-4}{}\fi
\noindent\textit{Source reference.} Original lemma 03JA: sources/stacks-chow/tags.txt, source lines 4632--4632. Cited at proof lines 1770-1773: provides the sufficiently divisible integer n.\par\medskip
\ifdefined\hypertarget\hypertarget{sourcepointer-5}{}\fi
\noindent\textit{Source reference.} Original lemma 0F3B: sources/stacks-chow/tags.txt, source lines 18861--18861. Cited at proof lines 1780-1782: keeps the identity w\_alpha = n w under base change in the reduction.\par\medskip
\ifdefined\hypertarget\hypertarget{sourcepointer-6}{}\fi
\noindent\textit{Source reference.} Original lemma more-morphisms-lemma-weighting-specialization: sources/stacks-more-morphisms/pinned-original-source.tex, source lines 22469--22469. Cited at proof lines 1795-1797: kills the difference weighting vanishing over the generic point of the discrete valuation ring.\par\medskip
\begin{lemma}
\label{lemma-weightings}
Let $S$ be a locally Noetherian scheme. Let $f : X \to S$ be a locally
quasi-finite morphism of schemes.
\begin{enumerate}
\item For $\alpha \in z(X/S, 0)$ the map $w_\alpha : X \to \mathbf{Z}$
constructed above is a weighting.
\item If $X$ is quasi-compact, then given a weighting
$w : X \to \mathbf{Z}$ there exists an integer $n > 0$ such
that $nw = w_\alpha$ for some $\alpha \in z(X/S, 0)$.
\item The integer $n$ in (2) may be chosen to be a power of
the prime $p$ if $S$ is a scheme over $\mathbf{F}_p$.
\end{enumerate}
\end{lemma}

\begin{proof}
Let $\alpha \in z(X/S, 0)$ and choose a diagram
$$
\xymatrix{
X \ar[d]_f & U \ar[l]^h \ar[d]^\pi \\
Y & V \ar[l]_g
}
$$
as in More on Morphisms, Definition \ref{more-morphisms-definition-weighting}.
Denote $\beta \in z(U/V, 0)$ the restriction of the base change $g^*\alpha$.
By the compatibility with base change (Lemma \ref{lemma-weightings-pre})
we have $w_\beta = w_\alpha \circ h$ and it suffices to show that
$\int_\pi w_\beta$ is locally constant on $V$. Next, note that
\begin{align*}
\left( \int_\pi w_\beta \right)(v)
& =
\sum\nolimits_{u \in U, \pi(u) = v} 
\beta(u) [\kappa(u) : \kappa(v)]_i [\kappa(u) : \kappa(v)]_s \\
& =
\sum\nolimits_{u \in U, \pi(u) = v} \beta(u)[\kappa(u) : \kappa(v)]
\end{align*}
This last expression is the coefficient of $v$ in $\pi_*\beta \in z(V/V, 0)$.
By Lemma \ref{lemma-relative-cycle-smooth} this function is locally
constant on $V$.

\medskip\noindent
Conversely, let $w : X \to S$ be a weighting and $X$ quasi-compact.
Choose a sufficiently divisible integer $n$. Let $\alpha$ be the family
of $0$-cycles on fibres of $X/S$ such that for $s \in S$ we have
$$
\alpha_s =
\sum\nolimits_{f(x) = s} \frac{n w(x)}{[\kappa(x) : \kappa(s)]_i} [x]
$$
as a zero cycle on $X_s$.
This makes sense since the fibres of $f$ are universally bounded
(Morphisms, Lemma
\ref{morphisms-lemma-locally-quasi-finite-qc-source-universally-bounded})
hence we can find $n$ such that the right hand side is an integer for
all $s \in S$. The final statement of the lemma also follows, provided
we show $\alpha$ is a relative $0$-cycle. To do this we have to
show that $\alpha$ is compatible with specializations along discrete
valuation rings. Note that neither the construction of $w_\alpha$
nor the proof of Lemma \ref{lemma-weightings-pre} uses that $\alpha$
is a relative cycle. Hence we have $w_\alpha = nw$ and this remains
true after base change. As well, the base change of a weighting
is a weighting, see
More on Morphisms, Lemma \ref{more-morphisms-lemma-weighting-base-change}.
Thus we reduce to the problem studied in the next paragraph.

\medskip\noindent
Assume $S$ is the spectrum of a discrete valuation ring with generic
point $\eta$ and closed point $0$. Let $w : X \to S$ be a weighting with
$X$ quasi-finite over $S$. Let $\alpha$ be the family of $0$-cycles on
fibres of $X/S$ constructed in the previous paragraph (for a suitable $n$).
We have to show that $sp_{X/S}(\alpha_\eta) = \alpha_0$.
Let $\beta \in z(X/S, 0)$ be the relative $0$-cycle on $X/S$
with $\beta_\eta = \alpha_\eta$ and $\beta_0 = sp_{X/S}(\alpha_\eta)$.
Then $w' = w_\beta - nw : X \to \mathbf{Z}$
is a weighting (using the result above) and zero in the points of $X$
which map to $\eta$. It follows from More on Morphisms,
Lemma \ref{more-morphisms-lemma-weighting-specialization}
that $w' = 0$. This means $w_\beta = nw$, hence $\alpha = \beta$ as desired.
\end{proof}

\end{document}
